
% 来源：EPUB part028.xhtml
\section{买股票就是买公司}

\BookQuestion{网友：投资大师的一句话也许就改变了我们的一生，望段总不吝赐教，感激不尽。}

如果有这么一句话，可能就是“买股票就是买公司”。我曾对巴菲特讲这是我从他那里学到的最重要的一句话，然后他说，这正是他从格雷厄姆那儿学到的最重要的一句话。（2010-05-23）

% 来源：EPUB part029.xhtml
\subsection{投资是什么？}

\BookMinorHeading{1. 基本版}

投资就是买未来现金流。

所谓能看懂公司就是能看懂其未来现金流（做对的事情）。

所有所谓有关投资的说法实际上都是在讨论如何看懂现金流的问题（如何把事情做对）。

% 来源：EPUB part030.xhtml
\BookMinorHeading{2. 说明版}

买股票就是买公司，买公司就是买其未来现金流折现。这里现金流指的是净现金流，未来指的是公司的整个生命周期。

折现率实际上是相对于投资人的机会成本而言的。最低的机会成本就是无风险回报率，比如美国国债的利率。

所谓能看懂公司就是能看懂其未来现金流（做对的事情）。

所有所谓有关投资的说法实际上都是在讨论如何看懂现金流的问题（如何把事情做对），比如生意模式、护城河、能力圈等。（2012-04-05）

% 来源：EPUB part031.xhtml
\BookMinorHeading{3.啰嗦版}

买股票就是买公司，买公司就是买其未来现金流折现。这里现金流指的是净现金流，未来指的是公司的整个生命周期，不是3年，也不是5年。

折现率实际上是相对于投资人的机会成本而言的。最低的机会成本就是无风险回报率，比如美国国债的利率。有些人把自己生意中有限的资金投到股市里实际上往往是不合算的，因为其自己的生意获利往往比股市的平均回报高。当然，多余资金投入无可非议。（听说国内某网络公司买了很多苹果的股票，这属于无可非议型的。）可我确实看到不少公司贷着款还要买股票，看不懂啊。

所谓能看懂公司就是能看懂其未来现金流折现（做对的事情）。所谓未来现金流折现只是个思维方式，千万不要去套公式，因为没人可以真的确定公式中的变量，所有假设可能都是不靠谱的。

个人观点：其实区分是不是“价值投资”的最重要、也许是唯一的一点，就是“投资者”是不是在买未来现金流折现。事实上我的确见到很多人买股票的理由很多时候都和未来现金流无关，但却和别的东西有关，比如市场怎么看，比如打新股一定赚钱，比如重组的概念，比如……呵呵，电视里那些个分析员天天在讲的那些东西。我有时会面带微笑看看CNBC的节目，那些主持人经常满嘴说着专业名词，但不知道为什么他们说了这么多年也不知道自己在说啥。

所有有关投资的说法实际上都是在讨论如何看懂现金流的问题（如何把事情做对），比如生意模式、护城河、能力圈，等等。在巴菲特这里我学到的最重要的东西就是生意模式。以前虽然也知道生意模式重要，但往往是和其他很多重要的东西混在一起看的。当年巴菲特特别提醒我，应该首先看生意模式，这几年下来慢慢觉得确实应该如此。

护城河实际上我觉得是生意模式的一部分，好的生意模式往往具有很宽的护城河。

好的生意模式往往是好的未来现金流的保障。

知道自己的能力圈有多大，往往比自己能力圈有多大要重要得多！

我觉得安全边际（margin of safety）实际上应该指的是能力圈而不仅仅是价格。

在自己能力圈内的生意自己往往容易懂得多，对别人的不确定性往往对自己是很确定的。比如当年我投网易时，市场不看好的原因是很多人觉得游戏这个市场不是很大。而我自己由于在这个行业里的时间很长，所以很确定这个市场非常大（但也不知道到底有多大，事实上最后的结果比我看到的还要大）。

不要轻易去“扩大”自己的能力圈。搞懂一个生意往往是需要很多年的，不要因为看到一两个概念就轻易跳进自己不熟悉的领域或地方，不然早晚会栽的。比如有的朋友跳进印度市场，有的朋友跳进日本市场。

我总是假设市场绝大多数情况下是非常聪明的，除非我发现市场确实错了。（这句话是针对“逆向操作”说的。逆向思维很重要，但逆向操作和随波逐流都是不可取的。最重要的是理性的独立思考能力。）

很多人说很难看懂未来现金流。其实绝大多数公司的未来现金流我也是看不懂的。看不懂的就不碰。一年两年或许更久的时间里总会有目标出现的。有些公司的生意模式很好，但股价有时候太贵，那就只能等了。好在这些年来一直如此，每隔些年就来个股灾，往往那时好公司也会跟着稀里哗啦的。

对于大多数不太了解生意的人而言，千万不要以为股市是个可以赚快钱的地方。长期来讲，股市上亏钱的人总是多过赚钱的人的。想赌运气的人还不如去买彩票，起码自己知道中的机会小，不会下重注。

也有人说股场就是赌场。事实上，对把股场当赌场的人们而言，股场确实就是赌场，常赌必输！

用我这个办法投资，一生可能会失去无数机会，但犯大错的机会也很少（但依然没办法避免犯错）。

我经常听见有人在讲哪只哪只股票赚了几倍的故事，可他们就是不说总的成绩，你懂的。

顺便感谢一下“自助餐先生”。凡是觉得我写的这点东西有帮助的人都应该多看看他老人家的东西，我能讲的他都讲过好多次了。

想到哪写到哪，主要是为了能随时提醒自己。

有人问如何避免以为看懂实际又错了的问题，个人观点：错误是不可避免的，但待在能力圈内以及专注和用功可以大幅度减少犯错的机会。

啥时候当你觉得简单版就足够了的时候，你大概就可以了。（2012-04-06）

我这里有很重要的一点就是认为大多数人其实不碰股票就是最好的投资，除非你认为自己确实有自己了解的好公司处在便宜的价钱。（2012-04-09）

我对巴菲特的认识非常简单，就是买股票就是买公司，从来没变过。（2021-07-28）

投资其实不需要啥理论的。买股票就是买公司其实属于定义，因为股票就是公司的一部分。（2024-01-24）

一般来讲我要投资一家公司时，主要考虑两个重要的点：

（1）这家公司能长期获利（足够的利润）吗？

（2）公司获得的利润如何给到股东？（2013-04-24）

\BookQuestion{网友：大道你已经财富自由了，推动你持续做投资的动力是什么？}

投资很有意思，是个非常好玩的游戏，我很喜欢（享受/enjoy）——很多年前我问过巴菲特先生这个问题，这是他当时的回答（大致意思）。（2023-04-23）

% 来源：EPUB part032.xhtml
\subsection{投资的信仰}

这里说的信仰不是形而上的东西。

我理解的投资归纳起来就是：买股票就是买公司，买公司就是买公司的未来现金流折现，句号！

这就是我说的所谓信仰的意思，我是从骨子里相信，不会因为任何影响而动摇。

很多人一讲到投资就会马上冒出很多大家都似懂非懂的术语，大概就是因为没有这个信仰。

其实投资就是价值投资的意思，不然投资投的是啥？

公司价值是什么？就是公司未来现金流的折现。

未来是多久？就是直到永远的意思。

永远是多久？就是到公司结束为止。

公司结束是什么意思？就是包括卖掉在内的所有可能。

记得马克思说过：价格围绕价值上下波动。买股票花的是价格，买的是价值。

其实我对投资的理解就这些，其他都是这几句话衍生出来的。

简单不？对，简单但绝不容易！不懂企业的人绝对没办法看懂未来现金流。所谓懂的人往往也只能看懂自己能力圈内的很少很少企业的未来现金流折现，而且还一定是毛估估滴。

学会简单其实就不简单——微博里看到的一句话，和我的简单但绝不容易的意思差不多。

未来现金流折现的公式？未来现金流折现只是一种思维方式，只有在自己能力圈范围内的公司，投资人才能毛估估看明白。没人真用公式的，至少胜者都是不用公式的。（说自己有公式还有参数的人们都是不懂甚至可能是骗人的，别理他们。）

很难吗？还好吧。我自己从事企业经营只有10多年，之后开始投资到现在也差不多10年。在这10年里有兴趣的企业中，我大致看懂了不到10家（被排除的不算，投错的不算），重手投了5家，差不多两年一家，没想象的那么难吧？

什么，你投的企业比我多多了？那你要小心了，一般而言，超出一定量以后，投的企业越多赚得越少。

怎么才能看懂企业或者叫看懂企业的未来现金流？

有句话叫八仙过海，各显神通。不过我喜欢巴菲特说的那个办法，从来没看过别人的，也没有兴趣看。我觉得真看懂了巴菲特的人大概是不会对别人的投资办法感兴趣的，觉得我武断的人请先看懂巴菲特哈。

说自己一半巴菲特一半索罗斯的意思实际上是说自己既不懂巴菲特也不懂索罗斯。（这句话和巴菲特说的15\%费雪85\%格雷厄姆的意思不同。）

投钱（我觉得这里不叫投资为好）的办法有很多，一般而言，知道得越多，赚得越少。不然大学教这个的就应该赚得最多。

巴菲特说生意模式最重要。我自己大概对这句话想了好几年了，还在继续想，但确实越来越感觉到巴菲特这个说法有道理。和巴菲特吃顿饭没白吃，这句话就已经值个100顿饭了吧？（其实远远不止哈。）

什么是生意模式？我也说不清楚，也许大学里的教授能给个文点的定义吧。但是，我觉得我大致知道那到底是什么意思。

如果你还不明白什么叫生意模式的话，那就看看巴菲特那几篇在几个大学的演讲。

还没懂？那看看喜诗糖果、可口可乐、比亚迪、苹果，看看航空公司，看看那些做太阳能晶片的“光伏企业”……还没懂？那重头来，接着再看。10到20年之内看不懂都可以用这个办法，就是“重头来，接着再看”。

20年以后还是看不懂怎么办？那20年以后再说了。（2012-06-24）

我从看到巴菲特的第一分钟起就开始相信巴菲特了，从来没动摇过。心中没有巴菲特的人是很难相信巴菲特的，而且谁说都没用。（2010-07-04）

其实每个人都有一颗投机的心，所以才需要信仰。我对信仰的理解就是“做对的事情”，或者说知道是“不对的事情”就别做了。人们热爱做“不对的事情”是因为这类事情往往有短期诱惑。（2012-05-22）

要相信的是：做对的事情最后会有好结果，或者说做错的事情最后会有大麻烦。至于把事情做对则是个过程，有时候快有时候慢，需要投资者耐心等待。（2013-02-03）

投资的信仰指的是：相信长期而言股市是称重机。对没有信仰的人而言，股市永远是投票器。（2013-05-17）

我说的信仰就是你骨子里真的相信的东西，是不会被某些事情动摇的，比如长期利润及净现金流好的公司股价早晚会跟上来等等。（2017-02-27）

\BookQuestion{网友：大道，您有信仰吗？}

当然有信仰，比如买股票就是买公司。

\BookQuestion{网友：我指的是宗教之类的。}

我明白你说的是什么，现在轮到你来明白我说啥了哈。

\BookQuestion{网友：您有支持您人生的信仰吗？}

我是有信仰的，但我没有宗教，我相信有道理的东西。以投资为例，我只要知道一个公司长远能赚钱，它的股价一定能跟上，所以就能忽略市场短期的波动。如果你没有这种信仰，就可能每天都去看今天有什么消息，股票会涨还是会掉。但如果你懂了一个公司，不管怎么涨跌，它还是那家公司。所以需要理解它的企业文化和生意模式。

\BookQuestion{网友：大道的“宗教”信仰就是做对的事情。}

是的，所谓信仰就是要做对的事情。（2022-01-23）

\BookQuestion{网友：我没有信仰，更糟糕的是我不知道如何改变。}

你必须相信。（You must believe.）——乌龟大师（2012-05-20）

\BookQuestion{网友：比尔·盖茨说巴菲特有很强烈的宗教信仰感，这让我挺意外的。我记得您说您对宗教的态度和巴菲特差不多，您了解的巴菲特是怎样的？}

我理解的“宗教”大概就是信仰的意思，就是说心里有些“必须相信”的东西。我的宗教态度确实和巴菲特一样，我们都信仰他对投资的理解——买股票就是买公司。买股票就是买公司这句话会说的人非常多，但我几乎没见过骨子里真的理解并相信这句话的——我知道的人里面大概不超过5个（包括巴菲特和芒格在内）——当然，这是因为我知道的人很少的缘故。（2013-04-11）

% 来源：EPUB part033.xhtml
\subsection{风险是第一考量}

记得以前翻过一本书叫《富爸爸，穷爸爸》，初学投资的人如果觉得巴菲特的东西一时不好懂，可以看看这本书。我对投资的基本理解好像和这本书差不多。

假设我有10000元闲钱，也就是一时半会用不着的钱，应该怎么办呢？我可能有以下的一些选择：

1. 在床底下挖个洞埋起来——20年后还是10000元；

2. 存在银行或买国债，假设平均利息是6\% ——20年后大概是32000多元；

3. 投资在标准普尔500指数（Standard \& Poor's 500 Index，简称S\&P

500），近100年的平均年回报大约是9\% ——20年后是56000元左右。

所以只要能找到任何年回报大于等于9\%的投资我就可以考虑投。

问题是，拿利息绝对无风险，而投资则有可能血本无归。

所以风险是决定是否投资的第一考量。

在自己懂的东西上投资，最重要的就是能看到风险在哪。

不同的人由于懂的东西不一样而投资在不同的投资标的上，是很正常的，没有孰优孰劣之分，但短期甚至长期回报有可能会不一样。如果对企业和投资都不了解的人，当S\&P500低的时候买指数也是个很好的投资办法。A股的情况我不了解。

其他如房地产投资的道理也一样。房地产价格是否贵可以看看《富爸爸，穷爸爸》，我记得里面说得挺好。

无论如何，我觉得人生最重要的投资是投资在教育上。

父母的教育，小时候的经历和苦难，小学、中学、大学和研究生学到的东西对我经营企业和现在的投资都起到了很大的作用。（2010-03-24）

投资用的是闲钱，不然就是投机了。（2014-07-20）

价值投资的风险不该比开车出门大，投机的风险不会比去赌场小。最危险的可能是一知半解的那种：比如号称自己是价值投资者，也知道价值投资要拿长线的道理（知其然不知其所以然），结果买错股票又坚持长线持有。（2011-03-01）

\BookQuestion{网友：先不谈风险，一个生意四年能回本，能不能投？}

先不谈风险的投资能叫投资么？（2012-07-19）

% 来源：EPUB part034.xhtml
\subsection{看懂生意是基本功}

投资的基本功其实就是看懂企业，看懂生意，无他。（2024-08-10）

只买生意，不预测股价。（2014-01-16）

这是投资的本质，不论你懂不懂，它都在那里。（2014-01-17）

从企业所有人的角度考虑是关键！（2014-10-20）

“如果你不愿意拥有一只股票十年，那就不要考虑拥有它十分钟。”（巴菲特）

这是所谓价值投资者决定买股票时要过的第一关。（2012-07-27）

\BookQuestion{网友：投资中最简单的事是什么？}

对投资，我想来想去，总觉得只有一样东西最简单，就是当你买一只股票时，你一定是认为你在买这家公司，你可能拿在手里10年、20年，有这种想法后就容易判断很多。（2010-04-27）

\BookQuestion{网友：净利润、现金流和自由现金流，哪个指标是我们应该重点关注的？}

如果你拥有一家公司的话，你就知道什么是你应该重点关注的了。如果你没有拥有一家公司的感觉，怎么说也都是没用的。（2011-05-24）

\BookQuestion{网友：我们到底该从哪下手学习未来现金流折现这种思维方式？}

大概只有一个办法，那就是保有拥有这家公司的心态。你要真有这种心态，其实不用学就会了，因为那都是真金白银啊。我发现虽然很多人说起来都是要用拥有公司的心态去买股票，但绝大多数实际上只是停留在口头上，一到关键时刻就会掉链子。（2012-02-11）

\BookQuestion{网友：有没有一只你认为不会卖掉的股票？}

每只股票我买的时候都是打算不卖的。（2019-08-21）

\BookQuestion{网友：巴菲特说过，我是个好投资家，因为我是一个企业家。我是个好企业家，因为我是投资家。您认为投资跟经营有什么区别吗？}

没觉得投资和自己经营企业有什么大区别。不同的地方是投资时是你认同的人在经营。（2019-04-06）
